\section{Introduction and Related Work}
Partial differential equations \citep[PDEs,][]{evans2022partial, strauss2007partial} form the foundation of scientific and engineering models \citep{ames2016nonlinear}, governing phenomena such as fluid flow \citep{naz2008comparison}, heat transport \citep{crank1947practical}, elasticity \citep{smith1990domain}, electromagnetics \citep{10.1098/rstl.1865.0008}, 
and the ability to accurately solve PDEs is thus a key enabler for scientific and engineering progress across a broad range of applications
\citep{zachmanoglou1986introduction}. Classical numerical solvers, such as including finite difference \citep{leveque2007finite}, finite element \citep{johnson2009numerical}, and spectral methods \citep{trefethen2000spectral}, provide reliable and well-understood approximations, but their computational cost grows rapidly with problem size, nonlinear dynamics, and parameter variability. As a result, simulations in high resolution or across large parameter spaces remain prohibitively expensive for real-time, uncertainty-aware, or many-query settings.

Recent efforts have explored machine learning techniques to accelerate PDE solution. Physics-informed neural networks~\citep[PINNs,][]{cai2021physics, cai2021physicsheat, raissi2019physics} embed the governing equations into a neural-network loss function and can learn solution fields in addition with observation data. Operator-learning models, such as Fourier Neural Operators \citep{li2020fourier, li2023fourier} and DeepONets \citep{lu2019deeponet, lu2021learning}, aim to learn mappings from forcing terms or parameters to full-solution fields and enable fast and flexible inference once trained. Neural surrogates, e.g., \citep{kochkov2021machine,sanchez-gonzalez20a-GNS}, replace expensive numerical methods with models learned from simulations 
and reduced-order models \citep{tripathy2018deep, white2019multiscale, eason2014adaptive} further accelerate computation by approximating fine-resolution solutions using coarse representations refined by learned corrections. However, these ML-based PDE solvers produce deterministic predictions \citep{huang2024diffusionpdegenerativepdesolvingpartial} and may require large datasets or costly hyperparameter tuning. Ensuring physical consistency and robustness across varying boundary conditions and parameter regimes remains challenging, motivating the exploration of probabilistic and generative approaches for forward PDE modeling.

In parallel, flow-based generative models based on diffusion \citep{song2021scorebasedgenerativemodelingstochastic} and stochastic interpolants \cite{JMLR:v26:23-1605} have emerged as state-of-the-art tools for sampling from complex, high-dimensional distributions. These models gradually transform noise into structured samples via stochastic differential equations, enabling diverse and uncertainty-aware generation in domains such as images \citep{rombach2022high, dhariwal2021diffusion}, audio \citep{kong2020diffwave, huang2023make, liu2023audioldm}, and molecular design \citep{guo2024diffusion, yim2024diffusion}. Their probabilistic formulation allows sampling from full posterior distributions rather than returning a single point estimate. These properties make diffusion models a promising candidate for scientific simulation, where multiple physically valid solutions may exist.

\subsection{Related Work}
Within the generative modeling literature, various methods have been applied to solving PDEs which can be broadly separated into two different categories: \emph{data driven} and \emph{physics-informed}. 

The data driven approach models the underlying relationship between inputs and outputs by discovering correlations and patterns directly from the data without explicitly prescribing physical laws. In the context of PDEs, many generative methods have attempted to use the data driven approach. For instance, variational Autoencoders and generative adversarial networks have been used extensively to model PDEs~\citep{gonzalez2018deep, creswell2018generative, farimani1709deep, mirza2014conditional}. Flow-based generative models have also found success recently in this area, with applications in fluid field prediction \citep{yang2023denoising, kohl2023benchmarking, Qiang2024-DDPM-Airfoil}, weather forecasting \citep{Price2025-GenCast},
3-dimensional turbulent flows \citep{molinaro2024generative}, modeling buoyancy \citep{li2025generative}.

Physics-informed approaches are common in the PINN literature \citep{cai2021physics, cai2021physicsheat, raissi2019physics} but have seen less common adoption by the generative modeling community.

\citet{lienen2306zero} utilize Dirichlet boundary conditions as guidance when sampling, but do not incorporate the full equations as part of a loss function, therefore it is still primarily a data driven approach. \citet{shysheya2024on,rozet2023score} use diffusion models trained on PDE data in combination with inference-time (``zero shot'') guidance for data assimilation, i.e., to condition the generation on sparse measurements. Closest to our work is
DiffusionPDE \citep{huang2024diffusionpdegenerativepdesolvingpartial} which similarly uses guidance to incorporate the governing equations into the likelihood during the sampling process.
%

The guidance methods uses in these framework are not specific to PDEs, and many methods have been developed for controlled generation, predominantly in computer vision~\citep[see,][for recent surveys]{daras2024surveydiffusionmodelsinverse, chung2025diffusionmodelsinverseproblems}. Motivated by the need for diversity and statistical consistency, several recent contribution incorporate Sequential Monte Carlo~\citep[SMC,][]{wu2023-TDS,cardoso2024monte,pmlr-v267-ekstrom-kelvinius25b, zhao2025generativediffusionposteriorsampling} and Markov Chain Monte Carlo \citep{Dang2026inferencetime, janati2025a-MixtureGDM, corenflos2024conditioning, Kalaivanan2025ESSFlow} in the guidance procedure. However, we have yet to see an adaptation of such methodology to physics-informed regularization of generative PDE solvers.

\subsection{Contributions}

We propose a new method for guidance of pretrained generative PDE solvers leveraging SMC for improved efficiency. 
This enables data assimilation as well as physics-informed regularization for solving constrained PDE systems.

However, while existing SMC-based guidance methods in computer vision and other domains \citep{wu2023-TDS, cardoso2024monte, stevens2025sequential, dou2024diffusion, pmlr-v267-ekstrom-kelvinius25b, zhao2025conditional} are primarily motivated by statistical consistency, we challenge this interpretation. While these methods indeed show improved empirical performance compared to their non-SMC counterparts, we argue that this is \emph{despite the fact that the SMC algorithms are highly degenerate when applied to conditional generative models}. In other words, we argue that existing SMC-based guidance methods for high-dimensional generative models are effective, not because they achieve a high effective sample size, but due to their inherent ``multiple try'' sampling nature. We elaborate on this interpretation in Section~\ref{sec:smc-ea}.

Based on this insight, our contributions are:
\begin{itemize}
    \item We develop an SMC-based guidance method for data assimilation and physics-informed regularization for solving constrained PDE systems.
    \item We propose to use a second order stochastic proposal augmented with PDE guidance (referred to as the SOSaG proposal) for better integration of the diffusion models.
    \item While the basic SMC-formulation is theoretically consistent, we show that it's incompatible with the efficient SOSaG proposal. We therefore propose a generalization of the SMC framework that trades statistical consistency for improved empirical performance, effectively resulting in a hybrid ``vanilla guidance'' and SMC approach. 
    \item We demonstrate that this new method outperforms other guided sampling approaches \citep{huang2024diffusionpdegenerativepdesolvingpartial} on multiple benchmark datasets, as well as two and three species interacting PDE systems. 
\end{itemize}

